\chapter{NKMath : vecteurs, matrices, angles}

Tout ce qui bouge à l'écran passe par NKMath. Une position, une taille, une
rotation, une couleur, un rectangle : ce sont ses types. Ce chapitre les couvre,
puis nomme les pièges du calcul à virgule flottante --- ceux qui font qu'un
programme juste sur le papier donne un résultat faux à l'exécution.

\section{Les vecteurs}

\begin{center}
\begin{tabular}{@{}lll@{}}
\toprule
\textbf{Suffixe} & \textbf{Type} & \textbf{Alias} \\
\midrule
\texttt{f} & \texttt{float32} & \texttt{NkVec2f}, \texttt{NkVec3f}, \texttt{NkVec4f} \\
\texttt{d} & \texttt{float64} & \texttt{NkVec2d}, \texttt{NkVec3d}, \texttt{NkVec4d} \\
\texttt{i} & \texttt{int32}   & \texttt{NkVec2i}, \texttt{NkVec3i}, \texttt{NkVec4i} \\
\texttt{u} & \texttt{uint32}  & \texttt{NkVec2u}, \texttt{NkVec3u}, \texttt{NkVec4u} \\
\bottomrule
\end{tabular}
\end{center}

\begin{nkcode}{Exemple pedagogique}
NkVec2f position(100.f, 50.f);
NkVec2f vitesse(1.f, 0.f);
position = position + vitesse * dt;      // les operateurs sont surcharges

const math::NkVec2u taille = fenetre.GetSize();   // une TAILLE, non signee
\end{nkcode}

\begin{nkretenir}
Le suffixe n'est pas décoratif. \texttt{NkVec2u} pour une taille d'écran ---
elle ne peut pas être négative. \texttt{NkVec2i} pour une position en pixels ---
elle peut l'être. \texttt{NkVec2f} pour une position dans le monde. Choisir le
mauvais suffixe compile et donne une soustraction qui repasse par le haut.
\end{nkretenir}

\section{Matrices et quaternions}

\texttt{NkMat2}, \texttt{NkMat3}, \texttt{NkMat4} --- et leurs variantes
\texttt{f} et \texttt{d}. Une matrice $4\times4$ transporte une transformation
complète : translation, rotation, échelle.

\texttt{NkQuat} représente une rotation sans les blocages de cardan que
produisent trois angles d'Euler enchaînés.

\begin{nkretenir}
\textbf{Décision d'architecture du dépôt} : les transformations se stockent en
\textbf{translation, rotation et échelle séparées}, jamais en matrice unique ---
et \textbf{jamais de cisaillement}. Une matrice compose les trois de façon
irréversible : on ne peut plus retrouver la rotation seule pour l'animer.
\end{nkretenir}

\section{Angles, aléatoire, géométrie}

\begin{center}
\begin{tabular}{@{}ll@{}}
\toprule
\textbf{En-tête} & \textbf{Contenu} \\
\midrule
\texttt{NkAngle.h} & angles typés --- degrés et radians ne se confondent plus \\
\texttt{NkEulerAngle.h} & les trois angles, quand on en a besoin \\
\texttt{NkRandom.h} & générateur reproductible à graine égale \\
\texttt{NkRectangle.h} & \texttt{NkRect2f}, \texttt{NkRect2i} \\
\texttt{NkSegment.h} & segments \\
\texttt{NkColor.h} & couleurs \\
\texttt{NkRange.h} & intervalles \\
\texttt{NkSIMD.h} & calcul vectoriel accéléré \\
\bottomrule
\end{tabular}
\end{center}

\begin{nkretenir}
\textbf{Un générateur reproductible à graine égale n'est pas un défaut : c'est
la condition d'un banc.} Les trois jeux du dépôt tirent leurs dés avec
\texttt{NkRandom} et une graine fixée --- c'est ce qui permet à un banc de
rejouer exactement la même partie et de vérifier qu'elle se termine.
\end{nkretenir}

\section{Les fonctions}

\begin{nkcode}{Kernel/Foundation/NKMath/src/NKMath/NkFunctions.h}
math::NkSin(x)      math::NkCos(x)      math::NkTan(x)
math::NkSqrt(x)     math::NkAbs(x)      math::NkPow(x, y)
math::NkMin(a, b)   math::NkMax(a, b)   math::NkClamp(v, lo, hi)
math::NkLerp(a, b, t)
math::NkRound(x)    math::NkFloor(x)    math::NkCeil(x)
\end{nkcode}

\begin{nkpiege}
\textbf{On emploie NKMath, jamais \texttt{<cmath>}.} \texttt{std::sqrt} et
\texttt{math::NkSqrt} donnent le même nombre aujourd'hui ; mais la règle du dépôt
est qu'aucun \texttt{std::} n'entre, et NKMath porte en plus les variantes SIMD
et les types du moteur.
\end{nkpiege}

\begin{nkretenir}
\texttt{NkClamp(v, lo, hi)} borne une valeur ; \texttt{NkLerp(a, b, t)}
interpole. Ces deux-là suffisent à écrire la plupart des animations : une valeur
qui va de \texttt{a} vers \texttt{b} en fonction d'un \texttt{t} entre 0 et 1,
bornée pour ne pas dépasser.
\end{nkretenir}

\section{Les constantes}

\texttt{NK\_PI\_F} pour les \texttt{float32}, \texttt{NK\_PI\_D} pour les
\texttt{float64}.

\begin{nkpiege}
\textbf{Il n'existe pas de \texttt{NK\_PI} tout court.} Le dépôt a payé cette
absence : un module avait défini sa propre macro locale, perdue lors d'une
extraction. Employez le suffixe qui correspond à votre type --- et si vous
cherchez \texttt{NK\_PI} et ne le trouvez pas, ce n'est pas que $\pi$ manque.
\end{nkpiege}

\section{La virgule flottante, et ce qu'elle ne sait pas faire}

\begin{nkpiege}
\textbf{\texttt{==} sur des flottants est presque toujours une faute.}
\texttt{0.1f + 0.2f} ne vaut pas \texttt{0.3f} : ces nombres n'ont pas de
représentation exacte en binaire. Comparez à une tolérance près :

\begin{nkcode}{Idiome du depot}
if (math::NkAbs(a - b) < 1e-4f) { /* egaux, a la tolerance pres */ }
\end{nkcode}

\textbf{Et la tolérance est un CHOIX, pas une constante universelle.} Le dépôt
emploie \texttt{1e-4f} pour fusionner des sommets de maillage : trop grande, deux
sommets distincts se confondent ; trop petite, deux sommets identiques restent
séparés. Le bon epsilon dépend de l'échelle de vos données.
\end{nkpiege}

\begin{nkpiege}
\textbf{L'addition flottante n'est pas associative.} $(a+b)+c$ peut différer de
$a+(b+c)$. C'est pourquoi deux exécutions d'un même calcul GPU peuvent donner des
résultats légèrement différents si l'ordre de réduction varie --- sans qu'aucun
code ait changé.

Conséquence pratique : \textbf{ne concluez pas d'un écart minuscule entre deux
mesures qu'un défaut existe}, avant d'avoir mesuré le bruit de la mesure
elle-même.
\end{nkpiege}

\begin{nkbilan}
Les types portent leur précision dans leur suffixe : \texttt{f}, \texttt{d},
\texttt{i}, \texttt{u} --- et le choix a des conséquences, pas seulement un
style. Les transformations se stockent en T + R + S séparées, jamais en matrice
unique. \texttt{NkClamp} et \texttt{NkLerp} suffisent à la plupart des
animations. On n'emploie pas \texttt{<cmath>}. Et sur les flottants : jamais
\texttt{==}, une tolérance choisie, et l'addition n'est pas associative.
\end{nkbilan}

\begin{nkquestions}
\begin{enumerate}
\item Pourquoi \texttt{NkVec2u} pour une taille d'écran et \texttt{NkVec2i} pour
      une position ?
\item Pourquoi le dépôt refuse-t-il de stocker une transformation en matrice
      unique ?
\item Existe-t-il un \texttt{NK\_PI} ? Que faire si l'on ne le trouve pas ?
\item Pourquoi \texttt{a == b} est-il dangereux sur des \texttt{float32} ?
\item En quoi un générateur aléatoire reproductible sert-il un banc ?
\end{enumerate}
\end{nkquestions}

\begin{nkpratique}
Écrivez une fonction \texttt{NkVec2f Osciller(float32 t, NkVec2f centre, float32
rayon)} qui rend un point tournant autour du centre.

Puis animez-la dans une application \texttt{NkCanvasApp} : le point doit décrire
un cercle, et sa vitesse doit être \emph{indépendante} de la cadence d'affichage.
Indice : c'est \texttt{deltaTime} qui décide, pas le nombre d'images.
\end{nkpratique}

\begin{nkdemo}
Prouvez que \texttt{0.1f + 0.2f != 0.3f}.

Affichez la différence avec dix décimales. Puis trouvez, par essais, la plus
petite tolérance qui déclare ces deux nombres égaux. Enfin, refaites l'essai avec
des \texttt{float64} : la tolérance nécessaire change-t-elle ? De combien ?

Concluez sur ce que « choisir un epsilon » veut dire.
\end{nkdemo}

\begin{nklogique}
Un objet se déplace de \texttt{vitesse * deltaTime} à chaque image.

Un joueur a 60 images par seconde, un autre 144. Après une seconde, ont-ils
parcouru la même distance ? Démontrez-le.

Puis supposez que le code s'écrive \texttt{position += vitesse} sans
\texttt{deltaTime} : que devient la réponse, et pourquoi ce défaut ne se voit-il
jamais sur la machine du développeur ?
\end{nklogique}
